\subsection{多项式的定义}

设 I 为一个集合。我们回顾（III，p.～452），A 上关于 I 的自由交换代数记为 \(\mathbf{A}[(\mathbf{X}_i)_{i\in I}]\) 或 \(\mathbf{A}[\mathbf{X}_i]_{i\in I}\) 。该代数的元素称为关于未定元 \(\mathbf{X}_i\) （或在未定元 \(\mathbf{X}_i\) 中）系数在 A 中的多项式。让我们回顾，未定元 \(\mathbf{X}_i\) 是 \(i\) 在 A 上关于 I 的自由交换代数中的典范像；有时用另一个符号表示该像很方便，例如 \(\mathbf{X}_i'\) , \(\mathbf{Y}_i\) , \(\mathbf{T}_i\) 等。这一约定常通过如下短语引入：「设 \(\mathbf{Y} = (\mathbf{Y}_i)_{i\in I}\) 为一族未定元」；此时所讨论的多项式代数记为 \(\mathbf{A}[\mathbf{Y}]\) 。当 \(\mathrm{I} = \{1,2,\dots,n\}\) 时，用 \(\mathbf{A}[\mathbf{X}_1,\mathbf{X}_2,\dots,\mathbf{X}_n]\) 代替 \(\mathbf{A}[(\mathbf{X}_i)_{i\in I}]\) 。

对于 \(\nu \in \mathbf{N}^{(I)}\) 我们令

\[
\mathbf {X} ^ {\nu} = \prod_ {i \in I} \mathbf {X} _ {i} ^ {\nu_ {i}}.
\]

则 \((\mathbf{X}^{\nu})_{\nu \in \mathbf{N}^{(I)}}\) 是 A-模 \(\mathbf{A}[(\mathbf{X}_i)_{i\in I}]\) 的一组基。这些 \(\mathbf{X}^{\nu}\) 称为未定元 \(\mathbf{X}_i\) 的单项式。对于 \(\nu = 0\) ，我们得到 \(\mathbf{A}[(\mathbf{X}_i)_{i\in I}]\) 的单位元。每个多项式 \(u\in \mathbf{A}[(\mathbf{X}_i)_{i\in I}]\) 可以唯一地写成如下形式

\[
u = \sum_ {\nu \in \mathbf {N} ^ {(I)}} \alpha_ {\nu} X ^ {\nu}
\]

，其中 \(\alpha_{\nu}\in A\) ，且这些 \(\alpha_{\nu}\) 除有限个外均为零； \(\alpha_{\nu}\) 称为 u 的系数； \(\alpha_{\nu}X^{\nu}\) 称为 u 的项（通常，元素 \(\alpha_{\nu}X^{\nu}\) 称为 \(X^{\nu}\) 中的项），特别地， \(\alpha_{0}X^{0}\) 的项（它与 \(\alpha_{0}\) 等同）称为 u 的常数项。当 \(\alpha_{\nu}=0\) 时，我们滥用语言说 u 不含 \(X^{\nu}\) 中的元素；特别地，当 \(\alpha_{0}=0\) 时，我们说u是一个无常数项的多项式（III，第453页）。1的任何标量倍称为常数多项式。

设 B 为交换环， \(\rho: \mathbf{A} \to \mathbf{B}\) 为环同态。我们把 \(\mathrm{B}[(\mathbf{X}_i)_{i \in I}]\) 通过 \(\rho\) 视为 A-代数。于是，映射 \(\sigma\) —— 即 \(\mathbf{A}[(\mathbf{X}_i)_{i \in I}]\) 到 \(\mathrm{B}[(\mathbf{X}_i)_{i \in I}]\) 中的、将 \(\sum \alpha_\nu \mathbf{X}^\nu\) 变换为 \(\sum \rho(\alpha_\nu) \mathbf{X}^\nu\) 的映射 —— 是 A-代数的一个同态；如果 \(u \in \mathbf{A}[(\mathbf{X}_i)_{i \in I}]\) ，我们有时用 \(^p u\) 表示 \(u\) 在这个同态下的像。 \(\mathbf{B} \otimes_{\mathbf{A}} \mathbf{A}[(\mathbf{X}_i)_{i \in I}]\) 到 \(\mathrm{B}[(\mathbf{X}_i)_{i \in I}]\) 中的、由 \(\sigma\) 典范地定义的同态，对每一个 \(i \in I\) 将 \(1 \otimes \mathbf{X}_i\) 变换为 \(\mathbf{X}_i\) ；这是 B-代数的一个同构（III，第449页）。

设 M 为以 \((e_{i})_{i\in I}\) 为基的自由 A-模。存在唯一的保单位同态 \(\varphi\) ，它把对称代数 \(\mathbf{S}(\mathbf{M})\) 映入代数 \(\mathbf{A}[(\mathbf{X}_{i})_{i\in I}]\) ，使得 \(\varphi(e_{i})=\mathbf{X}_{i}\) （对每个 \(i\in I\) ），并且这个同态是同构（III，第506页）。这个同构称为典范同构。它使我们能够将对称代数的某些性质应用于多项式代数。例如，设 \((\mathrm{I}_{\lambda})_{\lambda\in L}\) 是 I 的一个划分。设 \(\varphi_{\lambda}\) 为把 \(P_{\lambda}=\mathbf{A}[(\mathbf{X}_{i})_{i\in I_{\lambda}}]\) 映入 \(P=\mathbf{A}[(\mathbf{X}_{i})_{i\in I}]\) 的同态，它将 \(X_{i}\) （作为 \(P_{\lambda}\) 的元素）变换为 \(X_{i}\) （作为 P 的元素）。于是这些 \(\varphi_{\lambda}\) 定义了代数 \(\otimes_{\lambda\in L}P_{\lambda}\) 到代数 P 中的同态，

且此同态为同构（III，第503页，命题9）。

设 \(\mathbf{E}\) 是一个 \(\mathbf{A}\) -模，并令 \(\mathbf{E} \otimes_{\mathbf{A}} \mathbf{A}[(\mathbf{X}_i)_{i \in I}] = \mathbf{E}[(\mathbf{X}_i)_{i \in I}]\) 。 \(\mathbf{A}\) -模 \(\mathbf{E}[(\mathbf{X}_i)_{i \in I}]\) 的元素称为未定元 \(\mathbf{X}_i\) 中的、系数在 \(\mathbf{E}\) 中的多项式。这样的多项式可以唯一地写成 \(\sum_{\nu \in \mathbf{N}^{(I)}} e_\nu \otimes \mathbf{X}^\nu\) ，其中 \(e_\nu \in \mathbf{E}\) ，且这些 \(e_\nu\) 除有限个

下标外均为零；我们经常写 \(e_{\nu}X^{\nu}\) 以代替 \(e_{\nu}\otimes X^{\nu}\) 。

